% This file was converted to LaTeX by Writer2LaTeX ver. 1.9.9
% see http://writer2latex.sourceforge.net for more info
\documentclass{article}
\usepackage{calc,amsmath,amssymb,amsfonts}
\usepackage[LGR,T1]{fontenc}
\usepackage[greek,italian]{babel}
\usepackage[style=numeric,backend=biber]{biblatex}
\author{HP}
\date{2026-07-28}
\begin{document}
Prefazione

Ci sono progetti che nascono, crescono e si concludono nel giro di pochi mesi.

Altri, invece, rimangono silenziosamente con noi per gran parte della vita.

Questo libro appartiene alla seconda categoria.

L$\text{\textgreek{’}}$idea di realizzare un videogioco per Commodore 64 non è nata mentre scrivevo queste pagine. Non è
nata neppure quando ho iniziato a lavorare come programmatore. È molto più antica. Appartiene agli anni in cui un
computer a otto bit sembrava capace di aprire mondi sconosciuti e ogni nuova istruzione del BASIC dava
l$\text{\textgreek{’}}$impressione di avvicinarsi un poco di più a qualcosa che sembrava quasi irraggiungibile.

Appartiene anche a un insieme di suoni e di attese che il tempo non è riuscito a cancellare: il clic deciso
dell$\text{\textgreek{’}}$interruttore di alimentazione, il fruscio del televisore mentre cercava il canale giusto, il
rumore del Datassette che iniziava a girare. Poi l$\text{\textgreek{’}}$attesa, lunga e incerta, durante la quale
bastavano un nastro consumato, una regolazione imprecisa o un gesto involontario per costringere a ricominciare tutto
da capo.

Eppure non importava davvero.

Ogni tentativo diventava una scoperta. Ogni errore obbligava a ragionare. Ogni programma che finalmente funzionava
sembrava una piccola magia.

Come molti ragazzi della mia generazione, anch$\text{\textgreek{’}}$io sognavo di creare un videogioco.

Non una semplice prova di programmazione, ma un vero videogioco. Uno di quelli che osservavamo con meraviglia, cercando
di immaginare come fosse possibile far muovere un personaggio, rilevare una collisione, generare un effetto sonoro o
governare un intero mondo dentro una macchina che oggi definiremmo estremamente limitata.

Quel sogno, però, è rimasto incompiuto.

La vita ha preso altre strade. Lo studio, il lavoro, i progetti professionali e le responsabilità hanno progressivamente
sostituito quelle lunghe giornate trascorse davanti al Commodore 64. Eppure, senza che me ne rendessi conto, qualcosa
di quel periodo è rimasto sempre presente.

Non il desiderio di tornare indietro.

Non la nostalgia fine a se stessa.

Piuttosto la curiosità di capire, con gli strumenti e l$\text{\textgreek{’}}$esperienza maturati negli anni, come avrei
affrontato oggi quella stessa sfida.

È proprio da questa domanda che nasce questo libro.

Non dall$\text{\textgreek{’}}$ambizione di scrivere un manuale sul BASIC del Commodore 64 e nemmeno dal semplice
desiderio di raccontare alcuni ricordi personali. Entrambe queste componenti sono presenti, ma nessuna delle due
rappresenta, da sola, il motivo per cui queste pagine esistono.

Il vero obiettivo è un altro.

Vorrei accompagnare il lettore nella realizzazione di un videogioco così come avrebbe potuto nascere negli anni Ottanta:
un problema alla volta, una scelta alla volta, un errore alla volta. Senza scorciatoie e senza soluzioni calate
dall$\text{\textgreek{’}}$alto. Cercando di ricostruire non soltanto il codice, ma soprattutto il modo di ragionare che
rendeva possibile costruire software su una macchina con risorse tanto limitate.

Il Commodore 64, infatti, non chiedeva di essere semplicemente utilizzato. Invitava a fare domande. Spingeva a
sperimentare, a sbagliare, a osservare ciò che accadeva e a riprovare. Le sue limitazioni non erano soltanto ostacoli:
diventavano parte del problema e, spesso, anche parte della soluzione.

Ogni byte di memoria, ogni ciclo di elaborazione, ogni caratteristica dell$\text{\textgreek{’}}$hardware costringeva a
riflettere sulle conseguenze di una scelta. In questo senso il Commodore 64 è stato, per molti di noi, molto più di un
computer. È stato un primo laboratorio di informatica e, in qualche modo, anche un primo maestro. Ci ha insegnato che
programmare non significa semplicemente impartire istruzioni a una macchina, ma imparare a formulare problemi,
costruire ipotesi e verificare se una certa idea possa davvero funzionare.

Per questo motivo il libro non segue la struttura di un manuale tradizionale.

Gli argomenti non compaiono perché è arrivato il momento di spiegarli, ma perché il videogioco ne rende necessaria
l$\text{\textgreek{’}}$introduzione. Ogni nuova conoscenza nasce da un problema reale, ogni soluzione modifica il
progetto e ogni capitolo rappresenta un passo avanti nella costruzione del programma.

Non studieremo prima una teoria per cercarle successivamente un$\text{\textgreek{’}}$applicazione. Procederemo nella
direzione opposta. Avremo un obiettivo concreto e, ogni volta che lo sviluppo ci porrà davanti a un nuovo ostacolo, ci
fermeremo per comprenderlo. Sarà il progetto stesso a indicarci quali domande porre e quali strumenti imparare a usare.

Il lettore, quindi, non assisterà dall$\text{\textgreek{’}}$esterno alla realizzazione del videogioco.

Ne prenderà parte.

Capitolo dopo capitolo, vedrà il programma nascere, cambiare, fallire in alcuni tentativi e migliorare attraverso le
soluzioni successive. Potrà seguirne il percorso utilizzando un Commodore 64 reale oppure un emulatore moderno; potrà
ritrovare sensazioni già vissute o scoprire per la prima volta un modo di programmare appartenente a
un$\text{\textgreek{’}}$epoca diversa.

Accanto a questo percorso tecnico ne scorre un altro, più discreto.

È il racconto delle persone, degli incontri e delle esperienze che hanno acceso quella passione e che, molti anni dopo,
hanno reso possibile trasformarla in un$\text{\textgreek{’}}$opera compiuta.

La realizzazione del videogioco viene raccontata attraverso una storia vera non perché la mia vicenda personale
costituisca il centro del libro, ma perché è proprio da quella storia che nasce il progetto. Le persone incontrate, gli
insegnamenti ricevuti, gli entusiasmi condivisi, gli errori commessi e le occasioni impreviste rappresentano il
contesto umano dal quale prende forma il percorso tecnico che il lettore seguirà nelle pagine successive.

Narrazione e programmazione procedono quindi insieme.

Le due dimensioni non potrebbero essere del tutto separate, perché dietro ogni programma esistono sempre persone,
maestri, intuizioni, tentativi e domande che finiscono per influenzare il modo in cui impariamo a pensare.

Non so quale sarà il risultato di questo percorso per ciascun lettore.

Forse, durante la lettura, qualcuno ricorderà i pomeriggi trascorsi a copiare un listato da una rivista, la
soddisfazione del primo programma funzionante o l$\text{\textgreek{’}}$emozione di vedere qualcosa muoversi sullo
schermo per la prima volta. Qualcun altro incontrerà il Commodore 64 senza aver vissuto quegli anni e proverà a
comprendere perché una macchina così essenziale sia riuscita ad accendere la curiosità di
un$\text{\textgreek{’}}$intera generazione.

Mi auguro che, arrivato all$\text{\textgreek{’}}$ultima pagina, il lettore non abbia semplicemente imparato qualcosa in
più sul Commodore 64.

Spero che abbia avuto la sensazione di aver partecipato davvero alla nascita di un videogioco e, attraverso quella
costruzione, di aver imparato a guardare un computer con occhi diversi.

Perché, in fondo, è questa l$\text{\textgreek{’}}$esperienza che avrei voluto vivere quando tutto ebbe inizio.

E che, molti anni dopo, ho finalmente avuto l'occasione di raccontare attraverso queste pagine. Perché programmare
significa proprio questo: non domandarsi soltanto che cosa una macchina sappia fare, ma cominciare a chiedersi come sia
possibile insegnarglielo. È da domande come questa che nascono i programmatori.

Ed è da quella stessa domanda che, molti anni fa, è cominciato anche il mio viaggio.
\end{document}
